\subsection*{Logarithmic derivatives}
% \label{ostrich}

In Lie algebroid language, a Maurer--Cartan form on~$\Sigma $ is
nothing more than a morphism
$\omega \colon {T \Sigma } \rightarrow {\mathfrak g}$ of Lie
algebroids, and Theorem \ref{tslT} a special case of Lie II, as is
well-known. In~the general setting, we replace the Maurer--Cartan form
on $G$ with the action algebroid ${\mathfrak g} \times M$ asso\-ci\-a\-ted
with the action of $G$ on $M$, and use $f \colon \Sigma \rightarrow M$
to pull ${\mathfrak g} \times M$ back to a~Lie algebroid~$A(f)$ over
$\Sigma $. Of course this pullback must be performed {\em in the
 category of Lie algebroids} rather than vector bundles (see, e.g.,
\cite[Section~4.2]{Mackenzie_05}). The composite
$ \delta f \colon A(f) \rightarrow {\mathfrak g}$ of the natural maps
$A(f) \rightarrow {\mathfrak g} \times M \rightarrow {\mathfrak g} $
is a Lie algebroid morphism, which becomes the {\it logarithmic
 derivative} of $f$. The results to be described here show that
$\delta f$~--- or more precisely an appropriate equivalence class of
$\delta f$, see Section~\ref{infch}~--- is a complete invariant of~$f$.

\begin{Example}[logarithmic derivative of an embedding] Suppose
$\Sigma \subset M$ is a submanifold and $f \colon \Sigma \rightarrow
M$ the embedding. Then $A(f)$ is the subbundle of the trivial bundle
${\mathfrak g} \times \Sigma \rightarrow \Sigma$ consisting of all
pairs $(\xi,x)$ having the property that the integral curve on $M$
through $x \in \Sigma $ of the infinitesimal generator $\xi^\dagger$
of $\xi \in {\mathfrak g} $ is tangent to $\Sigma \subset M$ at
$x$. The anchor of $A(f)$ is $(\xi,x) \mapsto \xi^\dagger(x)$ and the
bracket well-defined by
\begin{equation*} [X,Y]=\nabla_{\#X}Y-\nabla_{\#Y}X +\{X,Y\}.
\end{equation*} Here $\nabla $ is the canonical flat connection on
${\mathfrak g} \times \Sigma$ and, viewing sections of ${\mathfrak g}
\times \Sigma $ as ${\mathfrak g} $-valued functions,
$\{X,Y\}(x):=[X(x),Y(x)]_{\mathfrak g} $. The logarithmic derivative
$\delta f$ is the composite $A(f)\hookrightarrow {\mathfrak g} \times
\Sigma \rightarrow {\mathfrak g}$.
\end{Example}

\subsection*{Generalized Maurer--Cartan forms}% \label{frog}

Again let $M$ be a smooth manifold on which some Lie group $G$ is
acting from the left transitively~--- what is hereafter referred to as
a {\it homogeneous $G$-space}. With this data fixed, our next task is
to introduce axioms for Lie algebroid morphisms
$\omega \colon A \rightarrow {\mathfrak g}$, where $A$ is a Lie
algebroid over some manifold $\Sigma$, modeled on local properties of
logarithmic derivatives $\delta f$ of smooth maps
$f \colon \Sigma \rightarrow M$.

To this end, observe that {\it logarithmic derivatives map Lie
algebroid isotropy algebras isomorphically onto isotropy algebras of
the action of $G$ on $M$}. Specifically, if we denote the kernel of an
anchor map $\#\colon A \rightarrow T \Sigma$ by $A^\circ $, then, for
any $x \in \Sigma$, $A(f )_x^\circ$ and ${\mathfrak g}_{f(x)}$ have the
same dimension, and
\begin{equation}
\delta f(A(f)_x^\circ) = {\mathfrak g}_{f(x)}.
 \label{uht}
\end{equation}
The following axioms, then, are no stronger than
properties already satisfied by logarithmic derivatives:
\begin{enumerate}\itemsep=0pt
\item[M1] $A$ is transitive.
\item[M2] For some point $x_0 \in \Sigma $, the restriction
$\omega \colon A^\circ_{x_0} \rightarrow {\mathfrak g}$ is injective.
\item[M3] For some (possibly different) point $x_0 \in \Sigma $,
there exists $m_0 \in M$ such that $\omega\big(A^\circ_{x_0}\big) \subset
{\mathfrak g}_{m_0}$, which will be written $x_0\xrightarrow{\omega}m_0$.
\end{enumerate}
%
Using the shorthand defined in M3, \eqref{uht} reads $x
\xrightarrow{\delta f} f(x)$. A Lie algebroid morphism $\omega \colon
A \rightarrow {\mathfrak g}$ is called a {\it generalized
Maurer--Cartan form} if it satisfies M1--M3.

Contrary to the group case, a Lie algebroid is not necessarily
integrable (the Lie algebroid of some Lie groupoid) and obstructions
to integrability are subtle. See \cite{Crainic_Fernandes_11} for a
fuller discussion and examples. Nevertheless, we have:
\begin{Proposition}\label{frogP}
Assume {\normalfont M1} and {\normalfont M2} hold
and that $\Sigma$ is connected. Then:
 \begin{enumerate}\itemsep=0pt
 \item[$1.$] $A$ is an integrable Lie algebroid.%\label{eins}
 \item[$2.$] {\normalfont M2} holds with $x_0$ replaced by an arbitrary
point $x \in \Sigma $.%\label{zwei}
 \item[$3.$] If {\normalfont M3} holds, then it holds with $x_0$
replaced by an arbitrary point $x \in \Sigma $ $($and suitable choice of
replacement $m \in M$ for $m_0)$.%\label{drei}
 \end{enumerate}
\end{Proposition}
%
We shall see the first assertion readily implies the others. A simple
proof of (1) is not known to us.\footnote{In an earlier
 version of this manuscript the main existence theorem was proven
 without assuming integrability, using substantially more complicated
 arguments, and integrability established post facto.} In Section
\ref{newsec}, where the proposition is proven, we will easily deduce
integrability from Crainic and Fernandes' generalization of Lie III
\cite{Crainic_Fernandes_03}.

\subsection*{Principal primitives}
% \label{pets}

In fact, the most na\"ive notion of a primitive is not unique ``up to
symmetry''. However, the na\"ive notion will play a role and be given a
name:
\begin{Definition}
%\label{petsD}
 A smooth map $f \colon \Sigma \rightarrow M$ is a {\it principal
 primitive} of a generalized Maurer--Cartan form
 $\omega \colon A \rightarrow {\mathfrak g} $ if there exists a Lie
 algebroid morphism $L \colon A \rightarrow A(f)$ such that the
 following diagram commutes:
 \begin{equation}
 \begin{tikzcd} A \arrow{d}{L} \arrow{r}{\omega } & {\mathfrak g}
 \\ A(f). \arrow[swap]{ru}{\delta f} &
 \end{tikzcd}
 \label{carpet:court}
 \end{equation}
\end{Definition}
 Note that we do not assume $L $ is an isomorphism, or even
that $A$ and $A(f)$ have the same rank.

\subsection*{Monodromy obstructions to the existence of
 primitives}
% \label{outline}

We now offer this paper's main construction, and formulate the
existence part of our results. Let~$M$ be a homogeneous $G$-space and
$\omega \colon A \rightarrow {\mathfrak g} $ an associated generalized
Maurer--Cartan form. We~are going to explicitly describe the
obstruction to the existence of a principal primitive
$f \colon \Sigma \rightarrow M $ of $\omega $, where $\Sigma $ is the
base of $A$. The most natural description is in terms of some abstract
transitive Lie groupoid ${\mathcal G}$ integrating $A$, whose
existence is guaranteed by Proposition \ref{frogP}, and which we may
take to be source-simply connected, on account of Lie I. In~the next
subsection we will offer a more concrete interpretation using a
generalization of Cartan's development along paths.

According to Lie II, $\omega $ is the derivative of a unique Lie
groupoid morphism
\begin{equation}\label{global}
 \Omega \colon\ {\mathcal G} \rightarrow G,
\end{equation}
which we call the {\it global form} of the monodromy of $\omega $,
being the analogue of \eqref{otter}. Continuing the analogy, we
choose $x_0 \in \Sigma $ and $m_0$ such that M2 and M3 hold, and
attempt to define a~principal primitive
$f \colon \Sigma \rightarrow M$, mapping $x_0$ to $m_0$, by
$f(x)=\Omega(p)\cdot m_0 $, where $p \in {\mathcal G} $ is any arrow
from $x$ to $x_0$. In~this ambition we are successful, so long as $f$
is well-defined, i.e., provided
\begin{gather}
\Omega({\mathcal G}_{x_0}) \subset G_{m_0}.\label{newt}
\end{gather}
Here ${\mathcal G}_{x_0} \subset {\mathcal G}$ denotes
the group of all arrows $p \in {\mathcal G} $ beginning and ending at
$x_0$, and $G_{m_0}$ the isotropy at $m_0 $ of the action of $G$ on
$M$.

Now the algebroid isotropy $A^\circ_{x_0}$ is the Lie algebra of
${\mathcal G}_{x_0} $ and, by our hypothesis M3,
$\omega\big(A^\circ_{x_0}\big) \allowbreak\subset {\mathfrak g}_{m_0}$. Therefore,
\begin{equation} \Omega\big({\mathcal G}_{x_0}^\circ\big) \subset G_{m_0},
\label{warbler}
\end{equation}
where ${\mathcal G}_{x_0}^\circ$ is the connected
component of ${\mathcal G}_{x_0}$. Moreover, as the transitive Lie
groupoid ${\mathcal G} $ has simply-connected source-fibres, there is
a natural exact sequence
\begin{equation*}
1 \rightarrow {\mathcal G}_{x_0}^\circ \rightarrow
{\mathcal G}_{x_0}\xrightarrow{\rho} \pi_1(\Sigma, x_0) \rightarrow 1.
%\label{shipout}
\end{equation*} From this and \eqref{warbler} we obtain a map
$\Omega_{x_0}^{m_0} \colon \pi_1(\Sigma ,x_0) \rightarrow M$
well-defined by
\begin{equation*} \Omega_{x_0}^{m_0}(\rho ( p))=\Omega(p) \cdot m_0.
\end{equation*}
We call this the {\it pointed form} of the
monodromy. By construction, our requirement~\eqref{newt} is equivalent
to $\Omega_{x_0}^{m_0}$ taking a constant value (which is necessarily~$m_0$).

\begin{Theorem}[existence and uniqueness of principal primitives]\label{outlineT}
 A Maurer--Cartan form $\omega{\rm :}$ $A \rightarrow {\mathfrak g} $
 over a connected manifold $\Sigma $ admits a principal primitive
 $f \colon \Sigma \rightarrow M$ if and only if the pointed form of
 the monodromy
 $\Omega_{x_0}^{m_0} \colon \pi_1(\Sigma, x_0) \rightarrow M$ is
 constant for some $($and consequently any$)$ choice of $x_0$ and $m_0$
 with $x_0\xrightarrow{\omega}m_0$. In~that case there is a unique
 principal primitive $f$ of~$\omega $ such that $f(x_0)=m_0$.
\end{Theorem}

\begin{proof}The preceding arguments establish the existence of a
primitive, given constant monodromy. Conversely, given the existence
of a primitive $f$ with $m_0=f(x_0)$, one easily establishes constancy
of the monodromy. For example, an elementary observation stated later
as Proposition \ref{uniquenessP} shows that
 \begin{equation}
 f(x) = \Omega (x_0)\cdot m_0\label{runningshot}
 \end{equation}
 for any arrow $p \in {\mathcal G} $ from $x_0$ to $x$ and so, in particular, $m_0=\Omega (p)\cdot m_0$ for any $p \in
{\mathcal G}_{x_0}$. Since~\eqref{runningshot} applies to any
principle primitive with $m_0=f(x_0)$, the last statement of the
theorem also holds.
\end{proof}

\subsection*{Monodromy as development along $\boldsymbol{A}$-paths}

Let $A$ be a Lie algebroid over a connected manifold $\Sigma $. Then a
piece-wise smooth map $a \colon [0,1] \allowbreak\rightarrow A$, covering an
ordinary path $\gamma \colon [0,1] \rightarrow \Sigma $, is called an
{\it $A$-path} if $\# a(t) = \dot \gamma(t)$, for all $t \in
[0,1]$. Here $\# \colon A \rightarrow T \Sigma $ denotes the anchor of
$A$.

Every $A$-path $a$ can be understood as Lie algebroid morphism $\hat a
\colon {T}I \rightarrow A$ defined by $\hat a (\partial/\partial t)=a$
and all such morphisms arise from $A$-paths. In~particular, given any
Lie groupoid~${\mathcal G} $ integrating $A$, we may apply Lie II,
obtaining a Lie groupoid morphism $I \times I \rightarrow {\mathcal G}
$. The image of $(0, t)$ under this morphism, denoted
\begin{equation*}
\integral_0^t a \in {\mathcal G},
\end{equation*} is an arrow from $\gamma(0)$ to $\gamma(t)$. If $A$ is
a Lie algebra, and ${\mathcal G} $ a Lie group with Lie algebra $A$, then
$a$ is simply a piece-wise smooth path in the Lie algebra, and the
integral above the usual integral to an element in the group. This
familiar case is the one applying in the proposition below:

\begin{Proposition}%\label{a:pathsP}
Consider a Maurer--Cartan form $\omega \colon A
 \rightarrow {\mathfrak g} $ as in the preceding theorem, and suppose
 $x_0\xrightarrow{\omega}m_0$, for some $x_0 \in \Sigma $ and
 $m_0 \in M$. Let $[\gamma] \in \pi_1(\Sigma,x_0)$ be given and let
 $a \colon [0,1]\rightarrow A$ be {\em any} $A$-path covering
 $\gamma $. Then the monodromy is given by
 \[
 \Omega_{x_0}^{m_0}([\gamma])=\bigg(\integral_0^1 \omega \circ a\bigg)\cdot m_0.
\]
\end{Proposition}
%
\begin{proof} The proposition follows immediately from the definition
of $\Omega_{x_0}^{m_0}$ and the following ele\-mentary property of
$A$-paths: Every Lie algebroid morphism $\omega \colon A_1 \rightarrow
A_2$ maps $A_1$-paths to $A_2$-paths, and if $\omega $ is the
derivative of a Lie groupoid morphism $\Omega \colon {\mathcal G}^1
\rightarrow {\mathcal G}^2$, then, for any $A_1$-path $a$,
 \begin{equation*}
 \Omega \bigg(\integral_0^t a\bigg) =
\integral_0^t \omega \circ a, \qquad t \in I.\tag*{\qed}
 \end{equation*}\renewcommand{\qed}{}
\end{proof}
